

Po przeprowadzeniu rygorystycznej analizy matematycznej w Rozdziale 2, w której wykazaliśmy, że z odpowiednio wysokim prawdopodobieństwem algorytmy oparte na próbkowaniu skompresowanym potrafią przełamać klątwę wymiarowości, niniejszy rozdział poświęcony jest weryfikacji tych rezultatów w praktyce.
Teoretyczne gwarancje, opierające się na własności ograniczonej izometrii (RIP), macierzowych nierównościach Chernoffa oraz twierdzeniu Wedina, dostarczają solidnych fundamentów, jednak to eksperymenty numeryczne pozwalają w pełni ocenić stabilność i wydajność zaproponowanych metod w rzeczywistych warunkach obliczeniowych.Celem tego rozdziału jest szczegółowe omówienie procesu implementacji Algorytmu \ref{alg:jednowymiarowy} (dla $k=1$) oraz Algorytmu \ref{alg:wielowymiarowy} (dla $k \ge 1$), a także zbadanie ich zachowania w zależności od doboru hiperparametrów modelu.
W szczególności interesować nas będzie wpływ liczby kierunków pomiarowych $m_\Phi$, liczby punktów próbkowania $m_\mathcal{X}$ oraz kroku różnicy skończonej $\epsilon$ na jakość zrekonstruowanej macierzy parametrów liniowych $A$.Ponadto, weryfikacji poddana zostanie stabilność algorytmów w obecności zakłóceń.
Jak zasygnalizowano w analizie teoretycznej, metody oparte na minimalizacji normy $l_1$ wykazują określoną odporność na szum, jednak istotne jest zbadanie zjawiska tzw. zwijania szumu (ang. \textit{noise folding}) w warunkach systematycznych błędów pomiarowych.
W ramach testów odtworzymy eksperymenty numeryczne zaproponowane przez Vybírala i współautorów, w tym badanie funkcji testowej w bardzo wysokim wymiarze ($d=1000$), której pomiary zostały celowo zaburzone szumem Gaussa o różnej wariancji. Struktura niniejszego rozdziału przedstawia się następująco:
\begin{itemize}
\item \textbf{Podrozdział 3.1: Środowisko obliczeniowe i narzędzia programistyczne} -- Opis wykorzystanych języków programowania, bibliotek do optymalizacji wypukłej (rozwiązywania problemów minimalizacji $l_1$) oraz procedur algebry liniowej (m.in. dekompozycji SVD).
\item \textbf{Podrozdział 3.2: Szczegóły implementacyjne algorytmów} -- Omówienie translacji teoretycznych kroków algorytmów na konkretne operacje numeryczne, w tym metody generowania wektorów losowych z odpowiednich miar sferycznych oraz aproksymacji pochodnych kierunkowych.
\item \textbf{Podrozdział 3.3: Weryfikacja dokładności rekonstrukcji (przypadek bezszumowy)} -- Testy działania algorytmów na wysoce wielowymiarowych funkcjach testowych w wyidealizowanych warunkach dokładnych pomiarów, potwierdzające ich wysoką skuteczność .
\item \textbf{Podrozdział 3.4: Doświadczenia numeryczne w obecności szumu} -- Analiza odporności algorytmów na zaburzenia wartości funkcji (np. szum Gaussa $\mathcal{N}(0,1)$ skalowany parametrem $\nu \in \{0.1, 0.01, 0.001\}$) i prezentacja map sukcesu rekonstrukcji aktywnych współrzędnych .
\end{itemize}
Przeprowadzone doświadczenia numeryczne pozwolą nie tylko potwierdzić rygorystyczne oszacowania z Rozdziału 2, ale również dostarczą praktycznych wskazówek dotyczących kalibracji algorytmów w rzeczywistych zastosowaniach analizy danych i uczenia maszynowego.